\documentclass[10pt]{article}

\usepackage[margin=1in]{geometry}
\usepackage{amsmath,amssymb,bm,mathtools}
\usepackage{algorithm}
\usepackage{algpseudocode}
\usepackage{float}
\usepackage{microtype}
\usepackage[hidelinks]{hyperref}

\newcommand{\Mink}{\langle\cdot,\cdot\rangle_M}
\newcommand{\HH}{\mathcal{H}}
\newcommand{\Exp}{\operatorname{Exp}}
\newcommand{\Log}{\operatorname{Log}}
\newcommand{\Proj}{\Pi}
\newcommand{\dist}{d_m}
\newcommand{\R}{\mathbb{R}}

\title{LorentzNet-H: Lorentz-Equivariant Mass-Shell\
Flow Matching for Jet Constituents}
\author{}
\date{}

\begin{document}
\maketitle

\begin{abstract}
We describe Arm H, a Lorentz-equivariant message-passing vector field trained by
Riemannian flow matching on a regulated mass shell. The network uses evolving
auxiliary Lorentz vectors for message passing, while its output field is built
intrinsically from mass-shell logarithmic maps and normalized tangent projections
of two condition-derived reference vectors. This document specifies the geometry,
network, conditional prior, coupling, objective, and sampling procedure.
\end{abstract}

\section{Geometry and notation}

Four-vectors are ordered as $p=(p^0,p^1,p^2,p^3)=(E,p_x,p_y,p_z)\in\R^{1,3}$
and use metric signature $(+,-,-,-)$,
\begin{equation}
 \langle p,q\rangle_M=p^0q^0-\sum_{a=1}^3p^aq^a,
 \qquad \|p\|_M^2\equiv p^2=\langle p,p\rangle_M.
\end{equation}
For a fixed radius (regulator mass) $m>0$, the forward hyperboloid is
\begin{equation}
 \HH_m=\{p\in\R^{1,3}:p^2=m^2,\ p^0>0\}.
\end{equation}
This is the radius-$m$ version of the unit hyperboloid used by
Lou et al.~\cite[Table~2]{lou2020neural}; their convention is obtained with
$x=p/m$ and $\langle x,y\rangle_L=-\langle p,q\rangle_M/m^2$.

The tangent space and induced Riemannian metric of this hyperboloid are
the standard expressions~\cite[Table~2]{lou2020neural}:
\begin{equation}
 T_p\HH_m=\{u:\langle p,u\rangle_M=0\},\qquad
 g_p(u,v)=-\langle u,v\rangle_M,\qquad
 \|u\|_p=\sqrt{-\langle u,u\rangle_M}.
\end{equation}
The corresponding ambient-to-tangent projection is~\cite[Table~2]{lou2020neural}
\begin{equation}
 \Proj_p(v)=v-\frac{\langle p,v\rangle_M}{m^2}p.
 \label{eq:projection}
\end{equation}

For $p,q\in\HH_m$, the hyperbolic distance is~\cite[Table~2]{lou2020neural}
\begin{equation}
 \gamma(p,q)=\frac{\langle p,q\rangle_M}{m^2},\qquad
 \dist(p,q)=m\operatorname{arcosh}\!\bigl(\gamma(p,q)\bigr).
\end{equation}
The associated logarithmic and exponential maps are~\cite[Table~2]{lou2020neural}
\begin{align}
 \Log_p(q)&=
 \frac{\operatorname{arcosh}(\gamma(p,q))}
 {\sinh(\operatorname{arcosh}(\gamma(p,q)))}
 \bigl(q-\gamma(p,q)p\bigr),
 \label{eq:log-map}\\
 \Exp_p(u)&=\cosh\!\left(\frac{\|u\|_p}{m}\right)p+
 m\sinh\!\left(\frac{\|u\|_p}{m}\right)\frac{u}{\|u\|_p}.
 \label{eq:exp-map}
\end{align}
Both expressions use their continuous values at $q=p$ and $u=0$.

\begin{center}
\fbox{\begin{minipage}{0.91\linewidth}
\small\textbf{Geometric interpretation.}
With $w_{p\to q}=q-\gamma(p,q)p$ and
$\|w_{p\to q}\|_p=m\sinh(\operatorname{arcosh}\gamma(p,q))$, the same map is
$\displaystyle \Log_p(q)=\dist(p,q)w_{p\to q}/\|w_{p\to q}\|_p$.
Thus the log map is the unit tangent direction from $p$ toward $q$, scaled by
the geodesic distance.
\end{minipage}}
\end{center}

A jet with $N$ real constituents is modelled as a point on the standard product
manifold.  Its product metric and componentwise maps are likewise standard
\cite{chen2023flow}:
\begin{equation}
 \mathcal{M}_N=\HH_m^N,\qquad
 G_X(U,V)=\sum_{i=1}^N g_{x_i}(u_i,v_i),
 \label{eq:product-metric}
\end{equation}
where $X=(x_1,\ldots,x_N)$ and $U=(u_1,\ldots,u_N)\in T_X\mathcal{M}_N$.
The product maps act componentwise:
\begin{equation}
 \Exp_X(U)=(\Exp_{x_1}(u_1),\ldots,\Exp_{x_N}(u_N)),\qquad
 \Log_X(Y)=(\Log_{x_1}(y_1),\ldots,\Log_{x_N}(y_N)).
 \label{eq:product-maps}
\end{equation}
The choice to use this product manifold for jet constituents, and the masking of
padded slots in implementation, are model-specific conventions.

\section{Conditional inputs}

The vector field receives the current physical state $X\in\mathcal{M}_N$, flow
time $t\in[0,1]$, and condition $c=[c_{\rm type},N,p_{T,J},m_J]$. It also receives
the ordered references $r_0=e_t=(1,0,0,0)$ and $r_1=p_J$. The latter is constructed
from the conditioned jet attributes,
\begin{equation}
 p_J=\frac{1}{s}
 \begin{pmatrix}
 m_{T,J}\cosh\eta_J\\
 p_{T,J}\cos\phi_J\\
 p_{T,J}\sin\phi_J\\
 m_{T,J}\sinh\eta_J
 \end{pmatrix},\qquad m_{T,J}=\sqrt{p_{T,J}^2+m_J^2},
 \label{eq:jet-reference}
\end{equation}
where $s$ is the dataset momentum scale. These vectors are condition-derived;
they are not formed by summing the target constituents.

\section{LorentzNet-H vector field}

The network maintains invariant scalar features $h_i^\ell\in\R^{d_h}$ and
auxiliary Lorentz-vector features $y_i^\ell\in\R^{1,3}$, with $d_h=96$ and
$L=6$ layers. The physical state $x_i$ remains fixed during a network evaluation.
Only $y_i^\ell$ evolves through message passing. The graph is complete over real
constituents, with $\mathcal{N}(i)=\{j:j\neq i\}$.

\begin{algorithm}[H]
\caption{LorentzNet-H vector field $V_\theta(X,t,c;r_0,r_1)$}
\label{alg:network}
\small
\begin{algorithmic}[1]
\Require $X=(x_1,\ldots,x_N)$, $t$, $c$, $(r_0,r_1)$
\Ensure $V_\theta(X,t,c;r_0,r_1)\in T_X\mathcal{M}_N$
\State $y_i^0\gets x_i$ and $\tau(t)\gets[t,\sin(\pi t),\cos(\pi t)]$ for all $i$
\State $h_i^0\gets\phi_{\rm seed}([\psi(\langle y_i^0,r_0\rangle_M),\psi(\langle y_i^0,r_1\rangle_M)])+\phi_c(c)-\phi_t(\tau(t))$
\For{$\ell=0,\ldots,L-1$}
  \State $\bar h_i^\ell\gets\operatorname{LayerNorm}(h_i^\ell)$
  \For{each $j\in\mathcal{N}(i)$}
    \State $\xi_{ij}^\ell\gets[\psi(\|y_i^\ell-y_j^\ell\|_M^2),\psi(\langle y_i^\ell,y_j^\ell\rangle_M)]$
    \State $M_{ij}^\ell\gets\phi_e^\ell(\bar h_i^\ell,\bar h_j^\ell,\xi_{ij}^\ell)$; \quad $\alpha_{ij}^\ell\gets\phi_x^\ell(M_{ij}^\ell)$
  \EndFor
  \State $h_i^{\ell+1}\gets h_i^\ell+\phi_h^\ell\!\left(\bar h_i^\ell,|\mathcal{N}(i)|^{-1/2}\sum_{j\in\mathcal{N}(i)}M_{ij}^\ell\right)$
  \State $y_i^{\ell+1}\gets y_i^\ell+\varepsilon_y|\mathcal{N}(i)|^{-1/2}\sum_{j\in\mathcal{N}(i)}\alpha_{ij}^\ell(y_i^\ell-y_j^\ell)$
\EndFor
\For{each $j\in\mathcal{N}(i)$}
  \State $\xi_{ij}^L\gets[\psi(\|y_i^L-y_j^L\|_M^2),\psi(\langle y_i^L,y_j^L\rangle_M)]$
  \State $a_{ij}\gets\phi_f(h_i^L,h_j^L,\xi_{ij}^L)$; \quad $q_{ij}\gets\Log_{x_i}(x_j)/(m+\dist(x_i,x_j))$
\EndFor
\State $v_i^{\rm part}\gets|\mathcal{N}(i)|^{-1/2}\sum_{j\in\mathcal{N}(i)}a_{ij}q_{ij}$
\State $(b_{i0},b_{i1})\gets\phi_r(h_i^L,[\psi(\langle y_i^L,r_0\rangle_M),\psi(\langle y_i^L,r_1\rangle_M)])$
\State $q_{ir}^{\rm ref}\gets\Proj_{x_i}(r_r)/(m+\|\Proj_{x_i}(r_r)\|_{x_i})$, for $r\in\{0,1\}$
\State $v_i^{\rm ref}\gets b_{i0}q_{i0}^{\rm ref}+b_{i1}q_{i1}^{\rm ref}$
\State $v_{\theta,i}\gets v_i^{\rm part}+v_i^{\rm ref}\in T_{x_i}\HH_m$
\State \Return $V_\theta=(v_{\theta,1},\ldots,v_{\theta,N})$
\end{algorithmic}
\end{algorithm}

The signed logarithmic preprocessing in Algorithm~\ref{alg:network} is
\begin{equation}
 \psi(z)=\operatorname{sgn}(z)\log(1+|z|).
 \label{eq:signed-log}
\end{equation}
All learned maps $\phi_\bullet$ act on scalar invariant inputs. With
$\sigma=\operatorname{SiLU}$, their widths are
\begin{align}
 \phi_c &: d_c\to d_h\to d_h,& \phi_t &: 3\to d_h\to d_h,&
 \phi_{\rm seed}&:2\to d_h\to d_h,\\
 \phi_e^\ell &: (2d_h+2)\to d_h\to d_h,&
 \phi_h^\ell &:2d_h\to2d_h\to d_h,\\
 \phi_x^\ell &:d_h\to d_h\to1,&
 \phi_f &: (2d_h+2)\to d_h\to1,&
 \phi_r &: (d_h+2)\to d_h\to2.
 \label{eq:mlp-widths}
\end{align}
Every hidden linear layer is followed by $\sigma$; $\phi_e^\ell$ also applies
$\sigma$ after its second linear layer. The final layers of $\phi_x^\ell$, $\phi_f$,
and $\phi_r$ have weights initialized from $\mathcal{N}(0,10^{-6})$ and zero bias.
The auxiliary-vector residual scale is $\varepsilon_y=10^{-2}$.

There is no message gate, softmax, absolute-value operation, or positivity constraint.
Thus the graph aggregation and the final particle readout are signed sums. The
particle term is intrinsically tangent because $q_{ij}\in T_{x_i}\HH_m$, and the
reference term is tangent by Eq.~\eqref{eq:projection}. Therefore their sum is
the model output directly in $T_{x_i}\HH_m$.

For $\Lambda\in SO^+(1,3)$, the network is jointly Lorentz equivariant,
\begin{equation}
 V_\theta(\Lambda^{\times N}X,t,c;\Lambda r_0,\Lambda r_1)
 =\Lambda^{\times N}V_\theta(X,t,c;r_0,r_1).
 \label{eq:equivariance}
\end{equation}
The statement follows because MLP inputs are Lorentz scalars, learned vector
coefficients are scalars, and tangent maps commute with Lorentz transformations.
Holding the references fixed in the laboratory frame supplies conditional symmetry
breaking through the beam/time and jet directions.

\section{Mass-shell Riemannian flow matching}

\subsection{Source and data distributions}

Observed massless constituent momenta $\bar x_i=(\|\bm p_i\|_2,\bm p_i)$ are
lifted to the shell by preserving spatial momentum,
\begin{equation}
 \iota_m(\bar x_i)=\left(\sqrt{\|\bm p_i\|_2^2+m^2},\bm p_i\right).
 \label{eq:shell-lift}
\end{equation}
Let $X_1\sim q_1(\cdot\mid c)$ denote a lifted target jet. The conditional source
distribution $q_0(\cdot\mid c)$ is an axis-aligned lognormal prior. It samples
\begin{align}
 \Delta\eta_i,\Delta\phi_i&\stackrel{\rm iid}{\sim}\mathcal{N}(0,0.15^2),&
 z_i&\stackrel{\rm iid}{\sim}\mathcal{N}(0,1),\\
 w_i&=\frac{e^{z_i}}{\sum_j e^{z_j}},&
 \eta_i&=\eta_J+\Delta\eta_i,\qquad\phi_i=\phi_J+\Delta\phi_i.
\end{align}
Writing
\begin{equation}
 R=\left\|\sum_iw_i(\cos\phi_i,\sin\phi_i)\right\|_2,
 \qquad p_{T,i}=w_i\frac{p_{T,J}}{R},
 \label{eq:prior-pt}
\end{equation}
the unlifted source four-vector is
\begin{equation}
 \bar x_{0,i}=\frac{1}{s}
 (p_{T,i}\cosh\eta_i,\ p_{T,i}\cos\phi_i,\ p_{T,i}\sin\phi_i,\ p_{T,i}\sinh\eta_i),
\end{equation}
followed by the lift in Eq.~\eqref{eq:shell-lift}.

\subsection{Online geodesic coupling}

Particle indices do not identify physical objects. For every fresh source draw,
prior constituents are paired with target constituents through the per-jet
Hungarian assignment
\begin{equation}
 \pi^\star=\arg\min_{\pi\in S_N}\sum_{i=1}^N\dist(x_{0,\pi(i)},x_{1,i})^2.
 \label{eq:icp}
\end{equation}
The source is permuted according to $\pi^\star$ before the RFM path is formed.
This coupling is recomputed online for each training batch and does not reuse a
fixed noise-to-data pairing.

\subsection{Conditional geodesic path and objective}

For a coupled pair $(X_0,X_1)$ and $t\sim\mathcal{U}[0,1]$, RFM uses the
product-manifold geodesic
\begin{equation}
 X_t=\Exp_{X_1}\!\left((1-t)\Log_{X_1}(X_0)\right)
 =\Exp_{X_0}\!\left(t\Log_{X_0}(X_1)\right).
 \label{eq:rfm-path}
\end{equation}
Its conditional velocity, tangent at $X_t$, is
\begin{equation}
U_t(X_t\mid X_1)=\frac{\Log_{X_t}(X_1)}{1-t}\in T_{X_t}\mathcal{M}_N.
\label{eq:target-field}
\end{equation}
Equation~\eqref{eq:target-field} is the geodesic specialization of the
distance-gradient target in Chen and Lipman~\cite[Eqs.~(13)--(15)]{chen2023flow},
not a separate choice of target field.  Let $d_{\mathcal{M}}$ denote the geodesic
distance of the product metric.  Along the conditional geodesic,
\begin{equation}
 d_{\mathcal{M}}(X_t,X_1)=(1-t)d_{\mathcal{M}}(X_0,X_1),\qquad
 \|\nabla_{X_t}d_{\mathcal{M}}(X_t,X_1)\|_{X_t}=1,
 \label{eq:geodesic-distance-identities}
\end{equation}
and the standard distance-gradient identity is
\begin{equation}
 \Log_{X_t}(X_1)=-d_{\mathcal{M}}(X_t,X_1)
 \nabla_{X_t}d_{\mathcal{M}}(X_t,X_1).
 \label{eq:log-distance-gradient}
\end{equation}
For $\kappa(t)=1-t$, the Chen--Lipman target therefore reduces as
\begin{align}
 U_t(X_t\mid X_1)
 &=-d_{\mathcal{M}}(X_0,X_1)\nabla_{X_t}d_{\mathcal{M}}(X_t,X_1) \notag\\
 &=\frac{\Log_{X_t}(X_1)}{1-t}.
 \label{eq:target-field-derivation}
\end{align}
The numerator vanishes linearly as $t\to1$; implementation terminates just before
$t=1$ and uses a protected denominator.

The model is trained by the product-manifold squared error
\begin{equation}
 \mathcal{L}_{\rm RFM}(\theta)=
 \mathbb{E}\left[
 \left\|V_\theta(X_t,t,c;r_0,r_1)-U_t(X_t\mid X_1)\right\|_{X_t}^2
 \right],
 \label{eq:rfm-loss}
\end{equation}
where the norm is defined by Eq.~\eqref{eq:product-metric} and is averaged over
real particles in implementation. Under squared-error regression, the population
optimum is the marginal velocity
\begin{equation}
 V^\star(x,t,c)=\mathbb{E}[U_t(X_t\mid X_1)\mid X_t=x,c].
 \label{eq:marginal-field}
\end{equation}
This field has the same time marginals as the mixture of conditional geodesic paths,
so integrating it transports $q_0(\cdot\mid c)$ toward $q_1(\cdot\mid c)$.

\subsection{Sampling}

Generation starts from $X(0)\sim q_0(\cdot\mid c)$ and integrates
\begin{equation}
 \dot X(t)=V_\theta(X(t),t,c;r_0,r_1).
 \label{eq:ode}
\end{equation}
Using geodesic Euler integration, one step is
\begin{equation}
 X_{k+1}=\Exp_{X_k}\!\left(\Delta t\,V_\theta(X_k,t_k,c;r_0,r_1)\right).
 \label{eq:euler}
\end{equation}
The product exponential is applied constituentwise. Hence all generated particles
remain on the positive-energy mass shell throughout sampling. The reported Arm H
configuration uses $m=0.1$, 64 Euler steps, and an endpoint time of $0.99999$.

\begin{thebibliography}{9}
\bibitem{lou2020neural}
Aaron Lou, Derek Lim, Isay Katsman, Leo Huang, Qingxuan Jiang, Ser-Nam Lim, and
Christopher M. De Sa.
\newblock Neural Manifold Ordinary Differential Equations.
\newblock In \emph{Advances in Neural Information Processing Systems}, 2020.
\newblock Appendix B.4, Table 2.

\bibitem{chen2023flow}
Ricky T. Q. Chen and Yaron Lipman.
\newblock Flow Matching on General Geometries.
\newblock \emph{arXiv preprint arXiv:2302.03660}, 2023.
\end{thebibliography}

\end{document}
